\chapter{Ciele}\label{ch:goals}

Hlavným cieľom práce je navrhnúť, implementovať a analyzovať dve architektonicky odlišné verzie tej istej webovej aplikácie – monolitickú a mikroservisnú – a skúmať ich správanie z hľadiska technických, prevádzkových a údržbových vlastností. Na dosiahnutie tohto cieľa sú stanovené nasledujúce podciele:
\begin{itemize}
    \item podrobný návrh a dokumentácia interného usporiadania monolitickej verzie aplikácie vrátane opisov modulov, vrstiev a dátových tokov;
    \item návrh mikroservisnej architektúry s jasným rozdelením doménových zodpovedností a definíciou API kontraktov;
    \item implementácia oboch verzií s použitím identického technologického stacku (React, Spring Boot~3, PostgreSQL, Docker Compose), aby bola zabezpečená objektívna porovnateľnosť;
    \item návrh a realizácia experimentov v kontrolovaných podmienkach na meranie odozvy systému, latencie pri interakcii s hernou doskou, spotreby systémových prostriedkov, priebehu spracovania požiadaviek a správania pri rôznych úrovniach záťaže;
    \item hodnotenie kvalitatívnych vlastností, ako sú modularita, čitateľnosť kódu, testovateľnosť, jednoduchosť údržby a operačná náročnosť;
    \item formulácia praktických odporúčaní, ktoré môžu pomôcť vývojárom pri rozhodovaní, kedy voliť monolit a kedy mikroservisnú architektúru.
\end{itemize}

Ciele sú definované tak, aby výsledok práce mohol slúžiť ako referenčná štúdia pre menšie a stredné projekty, kde rozhodnutia o architektúre výrazne ovplyvňujú náklady, škálovanie a dlhodobú udržateľnosť.